\PassOptionsToPackage{dvipsnames}{xcolor} % colori PDF/A

% setspace e footmisc vanno caricati PRIMA di pdfx (che carica hyperref): caricati dopo ridefiniscono \@footnotetext
% e i collegamenti delle note a piè di pagina perdono la destinazione ("Hfootnote.N has been referenced but does not exist").
% setspace: nessun suo comando è usato nel testo, ma insieme a \linespread rende compatti float e didascalie
% (toglierlo cambia l'impaginazione). footmisc[bottom]: le note stanno sempre in fondo alla pagina.
\usepackage{setspace}
\usepackage[bottom]{footmisc}

% xcolor va caricato prima di pdfx: se lo carica pdfx stesso, gli passa un'opzione (hyperref) ormai obsoleta e xcolor emette un avviso.
\usepackage[rgb]{xcolor}
% pdfx segnala sempre, con un avviso solo informativo, che i colori sono impostati in RGB (come richiede il PDF/A):
% il filtro nasconde solo questo messaggio.
\usepackage{silence}
\WarningFilter{pdfx}{Setting all color commands}
\usepackage{colorprofiles}
% PDF/A-1b (ISO 19005-1, livello di conformità B). È il livello che pdfLaTeX riesce davvero a soddisfare:
% il livello "a" (a-1a) richiede un PDF con tag di struttura, che pdfLaTeX non sa produrre, e dichiararlo senza
% esserlo fa fallire la validazione (regole 6.8.2.2 e 6.8.3.3). L'Ateneo richiede genericamente un "PDF/A".
% Con PDF/A-1 le immagini non possono avere trasparenza (PNG con canale alpha): usare JPEG o PNG senza alpha.
% Verifica: ./build.sh check (veraPDF) oppure https://www.pdf-online.com/osa/validate.aspx
\usepackage[a-1b,mathxmp]{pdfx}[2018/12/22]
% pdfx carica già hyperref (con l'opzione pdfa): non va ricaricato qui.
\usepackage[T1]{fontenc}
\usepackage[italian]{babel}
\usepackage{lmodern}
\usepackage{bookmark}
\usepackage{caption}
\usepackage{calc}
\usepackage{emptypage} % pagine vuote senza testatina e piè di pagina
\usepackage{indentfirst} % rientra il primo paragrafo di ogni sezione
\usepackage{graphicx} % immagini
\usepackage{longtable} % tabelle su più pagine
% glossaries carica datatool, che non ha il supporto per l'italiano: tracklang emetterebbe un avviso inutile
\usepackage{tracklang}
\ifdefined\TrackLangShowWarningsfalse\TrackLangShowWarningsfalse\fi
\usepackage[toc, acronym, automake]{glossaries}
\usepackage[backend=biber, style=verbose-ibid, hyperref, backref]{biblatex}
\usepackage[top=2.75cm, bottom=2.75cm, right=3cm, left=3.75cm]{geometry} % 1in+17pt+0.6cm
\usepackage{fancyhdr}
\usepackage{lipsum} % testo segnaposto: da togliere quando si scrive la tesi vera
\usepackage{titlesec}
% float fornisce la posizione [H] di figure e listati: prima arrivava per caso da minted.
% Va caricato qui, subito prima di minted: caricato prima cambia l'aspetto dei listati.
\usepackage{float}
\usepackage[cachedir=minted-caches]{minted} % https://it.overleaf.com/learn/latex/Code_Highlighting_with_minted
\usepackage{csquotes} % gestisce automaticamente i caratteri (")
\usepackage{etoolbox}

% Pacchetti facoltativi, non usati dal template: decommentare se servono.
% \usepackage{microtype}  % rifinitura tipografica (margini, spaziature): più lento da compilare e cambia l'impaginazione
% \usepackage{subcaption}  % sottofigure e sottotabelle (sostituisce il vecchio subfig)
% \usepackage{amsmath, amssymb, amsthm} % matematica avanzata
% \usepackage{booktabs}    % filetti professionali nelle tabelle
% \usepackage{comment}     % ambiente comment per escludere blocchi di testo
